\documentclass{article}
\usepackage[utf8]{inputenc}
\usepackage{geometry}
\usepackage{graphicx}
\usepackage{pgfplots}
\usepackage{amsmath}
\usepackage{amssymb}
\usepackage{amsfonts}
\usepackage[thmmarks, amsmath, thref]{ntheorem}
\usepackage{mdframed}
\usepackage{amsthm}
\usepackage{physics}
\usepackage{proof}
\usepackage{derivative}
\usepackage{comment}
\usepackage{hyperref}
\usepackage{wrapfig}
\usepackage{amsthm}
\usepackage{xcolor}
\renewcommand{\theenumi}{\roman{enumi}}

\theoremstyle{nonumberplain}
\theoremheaderfont{\itshape}
\theorembodyfont{\upshape}
\theoremseparator{.}
\theoremsymbol{\ensuremath{\square}}
\theoremsymbol{\ensuremath{\blacksquare}}
\newtheorem{solution}{Solution}
\theoremseparator{. ---}
\theoremsymbol{\mbox{\texttt{;o)}}}

\pgfplotsset{compat=1.18}
\begin{document}

El problema que tenemos es:

\[
\left\{
\begin{array}{l}
\dot{x} = 2t(x-1) \\
x(0) = 2
\end{array}
\right.
\]

Empecemos reescribiendo la ecuación diferencial:

\[
\dot{x} - 2t x = -2t
\]

Esta es una ecuación diferencial lineal de primer orden de la forma:

\[
\dot{x} + P(t)x = Q(t)
\]

donde \( P(t) = -2t \) y \( Q(t) = -2t \).


El factor integrante es:

\[
\mu(t) = e^{\int P(t)\,dt} = e^{\int -2t\,dt} = e^{-t^2}
\]


Multiplicamos ambos lados por \(\mu(t)\):

\[
e^{-t^2}\dot{x} - 2t e^{-t^2} x = -2t e^{-t^2}
\]

El lado izquierdo es la derivada del producto \(e^{-t^2}x\):

\[
\frac{d}{dt}\left(e^{-t^2} x\right) = -2t e^{-t^2}
\]


\[
\int \frac{d}{dt}\left(e^{-t^2} x\right) dt = \int -2t e^{-t^2} dt
\]

\[
e^{-t^2} x = \int -2t e^{-t^2} dt + C
\]

Observemos que \(\int -2t e^{-t^2} dt = e^{-t^2}\), ya que la derivada de \(-t^2\) con respecto a \(t\) es \(-2t\):

\[
e^{-t^2} x = e^{-t^2} + C
\]

 Despejamos \(x(t)\)

\[
x(t) = 1 + C e^{t^2}
\]

 Solo falta aplicar la condición inicial

Usamos \(x(0) = 2\):

\[
2 = 1 + C e^{0}
\implies C = 1
\]

Ordenando

\[
x(t) = 1 + e^{t^2}
\]

2. Método de sucesiones de Picard

Ahora, vamos a resolver usando el método de Picard y justificaremos los supuestos del Teorema de Picard-Lindelöf.

La ecuación integral equivalente es:

\[
x(t) = x_0 + \int_{0}^{t} 2s(x(s)-1) ds
\]
donde \(x_0 = 2\). \\

 Primero verificamos las hipótesis del Teorema de Picard-Lindelöf

La función \(f(t, x) = 2t(x-1)\):\\

- Es continua en \(t\) y \(x\) en todo \(\mathbb{R}^2\).\\
- Su derivada parcial respecto a \(x\) es \(\frac{\partial f}{\partial x} = 2t\), que es continua en \(t\).\\
- Por lo tanto, satisface la condición de Lipschitz local en \(x\).\\

Esto garantiza, por el Teorema de Picard-Lindelöf, la existencia y unicidad de la solución en algún intervalo alrededor de \(t=0\).\\

Vease ahora la construcción de la sucesión de Picard\\

Definimos la sucesión:

\[
x_{n+1}(t) = 2 + \int_{0}^{t} 2s(x_n(s)-1) ds
\]

Hacemos las primeras iteraciones:

Paso 0:
\[
x_0(t) = 2
\]

Paso 1:
\[
x_1(t) = 2 + \int_{0}^{t} 2s(2-1) ds = 2 + \int_{0}^{t} 2s\, ds = 2 + [s^2]_0^t = 2 + t^2
\]

Paso 2:
\[
x_2(t) = 2 + \int_{0}^{t} 2s(x_1(s)-1) ds = 2 + \int_{0}^{t} 2s((2 + s^2) - 1) ds = 2 + \int_{0}^{t} 2s(1 + s^2) ds
\]
\[
= 2 + \int_{0}^{t} 2s\, ds + \int_{0}^{t} 2s^3\, ds = 2 + [s^2]_0^t + [\frac{1}{2}s^4]_0^t = 2 + t^2 + \frac{1}{2}t^4
\]

Paso 3:
\[
x_3(t) = 2 + \int_{0}^{t} 2s(x_2(s)-1) ds = 2 + \int_{0}^{t} 2s((2 + s^2 + \frac{1}{2}s^4) - 1) ds
\]
\[
= 2 + \int_{0}^{t} 2s(1 + s^2 + \frac{1}{2}s^4) ds
\]
\[
= 2 + \int_{0}^{t} 2s\, ds + \int_{0}^{t} 2s^3\, ds + \int_{0}^{t} s^5\, ds
\]
\[
= 2 + [s^2]_0^t + [\frac{1}{2}s^4]_0^t + [\frac{1}{6}s^6]_0^t = 2 + t^2 + \frac{1}{2}t^4 + \frac{1}{6}t^6
\]

El proceso sigue n-ésimas veces, y se observa que la sucesión de Picard está generando la serie de Taylor de \(e^{t^2}\):

\[
x(t) = 1 + e^{t^2} = 1 + \sum_{n=0}^{\infty} \frac{t^{2n}}{n!}
\]

Por lo tanto, el método de Picard converge a la solución exacta encontrada previamente.
\begin{comment}
12. Suponga que \( F: U \subset \mathbb{R}^{n} \rightarrow \mathbb{R}^{n} \) es un campo vectorial y \( \Gamma \subset U \) una trayectoria simple y suave a tramos, parametrizada por \( \gamma:[a, b] \rightarrow \Gamma \subset \mathbb{R}^{n} \).\\
(a) Pruebe que
\[
\left|\int_{\Gamma} F \cdot \mathrm{~d} \gamma\right| \leq M\left(\ell_{\Gamma}\right)
\]

Recordemos la definición de la integral de línea
\[
\int_{\Gamma} F \cdot \mathrm{d}\gamma = \int_{a}^{b} F(\gamma(t)) \cdot \gamma'(t)\, dt
\]

Por la desigualdad del valor absoluto de una integral,
\[
\left| \int_{a}^{b} f(t) \, dt \right| \leq \int_{a}^{b} |f(t)| \, dt
\]
Aplicando esto a nuestra integral...

\[
\left| \int_{a}^{b} F(\gamma(t)) \cdot \gamma'(t)\, dt \right|
\leq \int_{a}^{b} \left| F(\gamma(t)) \cdot \gamma'(t) \right| dt
\]

Por la desigualdad de Cauchy-Schwarz:
\[
|F(\gamma(t)) \cdot \gamma'(t)| \leq \|F(\gamma(t))\| \cdot \|\gamma'(t)\|
\]
Así que:
\[
\left| \int_{a}^{b} F(\gamma(t)) \cdot \gamma'(t)\, dt \right|
\leq \int_{a}^{b} \|F(\gamma(t))\| \cdot \|\gamma'(t)\| dt
\]

Por definición, \( M = \max\{\|F(\bar{x})\| : \bar{x} \in \Gamma\} \), así que \( \|F(\gamma(t))\| \leq M \) para todo \( t \). Por lo tanto:
\[
\int_{a}^{b} \|F(\gamma(t))\| \cdot \|\gamma'(t)\| dt
\leq M \int_{a}^{b} \|\gamma'(t)\| dt
\]
Y
\[
\int_{a}^{b} \|\gamma'(t)\| dt = \ell_{\Gamma}
\]
es la longitud de arco de \( \Gamma \).

Unimos todo:
\[
\left|\int_{\Gamma} F \cdot \mathrm{d} \gamma\right| \leq M \ell_{\Gamma}
\]

donde \( M=\max \{\|F(\bar{x})\| \mid \bar{x} \in \Gamma\} \) y \( \ell_{\Gamma} \) es la longitud de arco de \( \Gamma \).


(b) Si en cada punto \( \gamma(t) \in \Gamma \) ocurre que \( F(\gamma(t)) \) esta en la misma dirección que el vector tangente a \( \Gamma \) en \( \gamma(t) \). Pruebe que
\[
\int_{\Gamma} F \cdot \mathrm{~d} \gamma=\int_{\Gamma}\|F\|\|\mathrm{d} \gamma\|
\]

Si \( F(\gamma(t)) \) es paralelo a \( \gamma'(t) \), entonces existe una función escalar \( \lambda(t) \geq 0 \) tal que:
\[
F(\gamma(t)) = \lambda(t) \gamma'(t)
\]
o, más generalmente, la dirección es la misma (podría ser negativa si permite sentido contrario):
\[
F(\gamma(t)) = \lambda(t) \frac{\gamma'(t)}{\|\gamma'(t)\|}
\]
Multipliquemos por \( \gamma'(t) \):
\[
F(\gamma(t)) \cdot \gamma'(t) = \lambda(t) \frac{\gamma'(t)}{\|\gamma'(t)\|} \cdot \gamma'(t) = \lambda(t) \|\gamma'(t)\|
\]
Pero también:
\[
\|F(\gamma(t))\| = |\lambda(t)|
\]
Por lo tanto, si \( \lambda(t) \geq 0 \), entonces:
\[
F(\gamma(t)) \cdot \gamma'(t) = \|F(\gamma(t))\| \cdot \|\gamma'(t)\|
\]
Así,
\[
\int_{\Gamma} F \cdot \mathrm{d} \gamma = \int_{a}^{b} F(\gamma(t)) \cdot \gamma'(t)\, dt = \int_{a}^{b} \|F(\gamma(t))\| \cdot \|\gamma'(t)\| dt
\]
Por definición,
\[
\int_{\Gamma} \|F\| \|\mathrm{d}\gamma\| = \int_{a}^{b} \|F(\gamma(t))\| \cdot \|\gamma'(t)\| dt
\]
por lo que
\[
\int_{\Gamma} F \cdot \mathrm{d} \gamma = \int_{\Gamma} \|F\| \|\mathrm{d}\gamma\|
\]
\end{comment}

\end{document}



